%%vsgtu709

%

% %%%

%%%   Самарский государственный технический университет

%%%   Кафедра прикладной математики и информатики

%%%

%%%   Преамбула для статьи в журнал "Вестник Самарского государственного

%%%   технического университета. Серия: «Физико-математические науки»"

%%%   Ориентирована под MikTEx 2.4 for Win32

%%%

%%%   Хотя сам журнал собирается под GNU/Linux на дистрибутиве TeXLive



\documentclass[11pt,a4paper,twoside]{article}



\usepackage[cp1251]{inputenc}

\usepackage[T2A]{fontenc}

\usepackage[english,russian]{babel}

\usepackage{samgtubib}

\usepackage[dvips]{graphicx}

\usepackage[multiple]{footmisc}



\usepackage{samgtu}

\renewcommand{\VestnikYear}{2009} % Год издания Вестника

\renewcommand{\VestnikNumber}{2\,(19)} % Номер Вестника

\renewcommand{\VestnikMonth}{Ноябрь} % Месяц выхода Вестника

\renewcommand{\VestnikData}{01.11.09} % Дата подписания Вестника в печать

\renewcommand{\VestnikUPL}{26,64} % Усл. печ. л.

\renewcommand{\VestnikUIL}{25,73} % Уч.-изд. л.

\renewcommand{\VestnikRegNumber}{479/09} % Регистрационный номер

\renewcommand{\VestnikOrder}{994} % Номер заказа







%%

% \begin{document}



%76S05	Flows in porous media; filtration; seepage

%86A05	Hydrology, hydrography, oceanography



%35K20	Second order parabolic equations, initial boundary value problems

%35Q53	KdV-like (Korteweg-de Vries) equations

%35G15	Boundary value problems for linear higher-order PDE



\renewcommand{\firstpage}{15}

\renewcommand{\lastpage}{24}

% \newcommand{\bigcom}{\text{\usefont{T2A}{cmfib}{m}{n}\large{\raisebox{1pt}{,}}}} % Cимвол производной в ТФКП (ковариантная)

% \Partition{Механика деформируемого твердого тела}{Applied mechanics} % Раздел Вестника, в который подаётся статья

% Названия разделов см. на сайте при подаче заявки





%%%%%%%%%%%%%%%%%%%%%%%%%%%%%%%%%%%%%%%%%%%%%%%%%%%%%%%%%%%%%%%%%%%%%%%%%%%%%%%%%%%%

%Информация о статье и об авторах на русском и английском языках



% Номер УДК

\udk{517.954}%Обработка текста. Подготовка текстов : Языки программирования. Компьютерные языки

\msc{35K70; 35K35, 35C15, 35A20}% Коды MSC см. http://www.ams.org/msc/





\title{Метод Римана для решения нелокальных краевых задач для псевдопараболических уравнений третьего порядка

}% Полный заголовок

{Метод Римана для решения нелокальных краевых задач для псевдопараболических уравнений \dots}

% Заголовок, который идёт в верхний колонтитул



\author{M.~Х.~Бештоков}% Автор(ы) для заголовка, если авторов несколько укать \inst

{М.~Х.~Б\,е\,ш\,т\,о\,к\,о\,в}% Автор(ы) для оглавления и колонтитула через \,



\institute{\kbgu.}

% Если несколько, то так  \institute{ Организация 1 \and Организация 2  \and Организация 3}



\email{E-mail}{beshtokov\_murat@rambler.ru} %E-mail's авторов



\otherinf{\fullauthor{Мурат Хамидбиевич Бештоков}\regalia{к.ф.-м.н.,

доц.}\position{докторант} \dept{каф. вычислительной математики} }



\keywords{краевые задачи, метод функции Римана, нелокальное условие,

уравнение в частных производных  третьего порядка,

псевдопараболическое уравнение}



\workabstract{% Здесь должна быть краткая аннотация статьи, например такая:

С помощью метода функции Римана доказаны

существование и единственность регулярных решений нелокальных

краевых задач для псевдопараболических уравнений третьего порядка с

переменными коэффициентами.}



\engtitle{Riemann method for solving non-local boundary value

problems for the third order pseudoparabolic equations}





\engauthor{M.~H.~Beshtokov}{M.~H.~B\,e\,s\,h\,t\,o\,k\,o\,v}





\enginstitute{\engkbgu.}



\otherenginf{\fullauthor{Murat H. Beshtokov} \regalia{Ph.\,D. Phys. \& Math.} \position{Doctoral Candidate} \dept{Dept. of~Computational Mathematics}}



\engabstract{The existence and uniqueness of regular solutions of non-local boundary value problems for the third order pseudoparabolic equations with variable coefficients are proved using the Riemann function method.}



\engkeywords{boundary value problems, the Riemann function method, non-local condition, partial differential equation of the third order, pseudoparabolic equation}



\date{10/I/2013} % Дата поступления статьи в редакцию.

\revisiondate{27/III/2013} % В окончательном виде



%%%%%%%%%%%%%%%%%%%%%%%%%%%%%%%%%%%%%%%%%%%%%%%%%%%%%%%%%%%%%%%%%%%%%%%%%%%%%%%%%%%%





\maketitle





Хорошо известно, что рассмотрение вопросов фильтрации жидкости в пористых средах \cite{besh:bib:Barenblat, besh:bib:Dzekzer}, передачи тепла в гетерогенной среде \cite{besh:bib:Rubinshtein, besh:bib:TingCooling}, влагопереноса в почво"=грунтах \cite{besh:bib:Hallaire} (см. \cite[c. 137]{besh:bib:Chudnovsky}) приводит к~модифицированным уравнениям диффузии, которые являются псевдопараболическими уравнениями в~частных производных третьего порядка.



В данной работе рассматриваются нелокальные краевые задачи для псевдопараболических уравнений третьего порядка с переменными коэффициентами. 

Доказательство существования регулярных решений поставленных задач проводится методом функции Римана, причём функция Римана, которая вводится здесь, существенно отличается от рассмотренной в \cite{besh:bib:Colton, besh:bib:Colton2}. 

В~работе \cite{besh:bib:Shhanukov} рассматриваются локальные и~нелокальные краевые задачи для псевдопараболических уравнений третьего порядка с переменными коэффициентами. 

Доказательство основных утверждений в \cite{besh:bib:Shhanukov} проводится с помощью метода функции Римана, принцип построения которой используется \mbox{в данной работе.}



\smallskip



\Section[n]{Задача А. Существование и единственность решения задачи А}

В замкнутом прямоугольнике $\overline{D}=\{(x,t): 0 \leq x \leq l,

0\leq t\leq T\}$ рассмотрим следующую нелокальную краевую задачу:

\begin{multline}

Lu=(\eta(x,t) u_{xt})_{x} +(k(x,t)u_{x})_{x}+ r(x,t)u_{x}+du_{t}-q(x,t)u =-f(x,t),

\\

0 < x < l,\ 0<t<T,

\label{besh:1}

\end{multline}

\begin{gather}

u(0,t)=\beta(t)\int^{l}_{0}u(x,t)dx +

\int^{t}_{0}\rho(t,\tau)u(l,\tau)d\tau-\mu(t), \quad 0 \leq t \leq T,

\label{besh:2}

\\

u_{x}(l,t)=0, \quad 0 \leq t\leq T, 

\label{besh:3}

\\

u(x, 0)=u_{0}(x), \quad 0 \leq x\leq l, 

\label{besh:4}

\end{gather}

где

\begin{equation}

\eta(x,t)\geq c_{0}>0, \ c_{0}={\rm const;}  \

\eta_{x,t},\ k_{x}, \ r_{x}, \ d_{t}, \ q,\ f \in C[\overline{D}];\

u_{0}(x) \in C^{2}[0, l], 

\label{besh:5}

\end{equation}

$\beta(t)$, $\rho(t,\tau)$, $\mu(t)$ "--- функции, непрерывные на $[0,T]$, $0\leq\tau\leq t$. Кроме этого, введём обозначение $\Pi(x,t)=ku_{x}+\eta u_{xt}$, которое будет использоваться далее.



Заметим, что нелокальное условие \eqref{besh:2} можно заменить условием

$$

u(0,t)=\beta(t)\int^{h}_{0}u(x,t)dx +

\int^{t}_{0}\rho(t,\tau)u(l,\tau)d\tau-\mu(t),

$$

где $h$ "--- глубина корнеобитаемого слоя (см. \cite{besh:bib:Chudnovsky2}) или активный слой почвы, который участвует в водоснабжении корневой системы, в~процессах испарения и транспирации. 

Поставленные и~исследованные в~данной работе задачи характерны также тем, что содержат в краевых

условиях нелокальность по времени, впервые изученную А.~И.~Кожановым \cite{besh:bib:Kozhanov}.







\smallskip



\begin{theorem}[1]Пусть коэффициенты уравнения \eqref{besh:1} и~граничных условий \eqref{besh:2}--\eqref{besh:4} удовлетворяют условиям гладкости \eqref{besh:5}{\rm ,}

$d(x, t)<0$ для любого $(x, t) \in D$ и $\beta(t)<0$ для любого $t \in [0, T].$ 

Тогда задача \eqref{besh:1}--\eqref{besh:4} имеет единственное регулярное в~$\overline{D}$ решение{\rm .}

\end{theorem}



\smallskip



\begin{proof}

Следуя \cite{besh:bib:Shhanukov} введём аналог функции Римана $w=w(x, t; \alpha, \tau)$ для уравнения \eqref{besh:1} в области $\Omega = \{(x,t) : \alpha<x<l, 0<t<\tau\}$ в~форме

\begin{multline}

Mw(x,t;\alpha,\tau)= \\ =-(\eta(x,t) w_{x})_{xt}+(k(x,t)w_{x})_{x}

-(r(x,t)w)_{x}

-(dw)_{t}- q(x,t)w=0,

\label{besh:6}

\end{multline}

\begin{gather*}

w(\alpha,t;\alpha,\tau)=0,

\quad 

w_{x}(\alpha,t;\alpha,\tau)={\eta^{-1}(\alpha,\tau)}

\exp \biggl({\int^{t}_{\tau}\frac{k(\alpha,t_{1})}{\eta(\alpha,t_{1})}dt_{1}}\biggr),

\\ 

w(x,\tau;\alpha,\tau)=\omega_{1}(x,\tau),

\end{gather*}

где $\omega_{1}(x,\tau)$ "--- решение следующей задачи Коши:

\begin{gather*}

\bigl(\eta(x,\tau)w_{x}(x,\tau;\alpha,\tau)\bigr)_{x}

+d(x,\tau)w(x,\tau;\alpha,\tau)=0,

\\

w(\alpha,\tau;\alpha,\tau)=0, \quad 

w_{x}(\alpha,\tau;\alpha,\tau)={\eta^{-1}(\alpha,\tau)}.

\end{gather*}

Здесь и далее $( \alpha, \tau)$ "--- произвольная фиксированная точка области $D$.



Имеет место соотношение

\begin{equation}

w Lu - uMw = \frac{\partial Q}{\partial x} - \frac{\partial

P}{\partial x},

\label{besh:7}

\end{equation}

где

\begin{gather*}

Lu=(\eta u_{xt})_{x}+(ku_{x})_{x} +ru_{x}+ du_{t} - qu + f(x,t),

\\

Q =\eta w u_{xt}+u(\eta w_{x})_{t} +kw u_{x}-k w_{x}u+ru w,

\quad 

P=\eta w_{x} u_{x}-duw.

\end{gather*}

Потребуем непрерывности $P$, $Q$ в $\overline{D}$, непрерывности и

ограниченности $P_{x}$, $Q_{t}$ в~$\overline{D}$.



Проинтегрируем соотношение \eqref{besh:7} по области $\Omega$:

\begin{equation}

\int_{\Omega}\big(w Lu - uMw\big)dxdt =

\int^{l}_{\alpha}\int^{\tau}_{0}\Big(\frac{\partial Q}{\partial x} -

\frac{\partial P}{\partial x}\Big)dxdt.

\label{besh:8}

\end{equation}

Тогда из \eqref{besh:8} с учётом определения функции $w=w(x,t; \alpha,\tau)$

получим представление

\begin{multline}

u(\alpha,\tau) = u(l,\tau)\eta(l,\tau)w_{x}(l,\tau;\alpha,\tau) -

\\

-\int^{\tau}_0\Big(\big(\eta(l,t)u_{xt}(l,t)+

k(l,t)u_{x}(l,t)\big)w(l,t;\alpha,\tau)+

\\

+u(l,t) \Big[\big(\eta(l,t)w_{x}(l,t;\alpha,\tau)\big)_{t}-

k(l,t)w_{x}(l,t;\alpha,\tau)+

r(l,t)w(l,t;\alpha,\tau)\Big]\Big)dt+

\\

+\int^{l}_\alpha\Big(d(x, 0)w(x, 0;\alpha,\tau)u(x, 0)-\eta(x, 0)w_{x}(x, 0;\alpha,\tau)u_{x}(x, 0)

\Big)dx-

\\

-\int^{\tau}_0 \int^l_{\alpha}w(x,t;\alpha,\tau)f(x,t)dxdt.

\label{besh:9}

\end{multline}

%где $(\alpha,\tau)$ "--- произвольная фиксированная точка области $D$.

Существование и единственность аналога функции Римана $w=w(x,t;\alpha,\tau)$ доказаны в работе~\cite{besh:bib:Shhanukov}.



Из представления \eqref{besh:9} простым преобразованием получим

\begin{multline}

u(\alpha,\tau) = u(l,\tau)\eta(l,\tau)w_{x}(l,\tau;\alpha,\tau)

-u_{x}(l,\tau)\eta(l,\tau)w(l,\tau;\alpha,\tau) + 

\\

+\int^{\tau}_0\Big(

u_{x}(l,t)\big[\big(\eta(l,t)w(l,t;\alpha,\tau)\big)_{t}-

k(l,t)w(l,t;\alpha,\tau)\big]-

\\

-u(l,t)

\Big[\big(\eta(l,t)w_{x}(l,t;\alpha,\tau)\big)_{t}

-k(l,t)w_{x}(l,t;\alpha,\tau)+ 

r(l,t)w(l,t;\alpha,\tau)\Big]\Big)dt+

\\

+\int^{l}_\alpha\Big(d(x, 0)

w(x, 0;\alpha,\tau)u(x, 0)-\eta(x, 0)w_{x}(x, 0;\alpha,\tau)u_{x}(x, 0)

\Big)dx-

\\

-\int^{\tau}_0 \int^l_{\alpha}w(x,t;\alpha,\tau)f(x,t)dxdt+

u_{x}(l, 0)\eta(l, 0)w(l, 0;\alpha,\tau).

\label{besh:10}

\end{multline}

Проинтегрируем \eqref{besh:10} по $\alpha$ от 0 до $l$. Тогда с учётом \eqref{besh:4} из

\eqref{besh:10} получим

\begin{multline}

\int^{l}_{0}u(\alpha,\tau)d\alpha =

u(l,\tau)\eta(l,\tau)\int^{l}_{0}w_{x}(l,\tau;\alpha,\tau)d\alpha

-u_{x}(l,\tau)\eta(l,\tau)\int^{l}_{0}w(l,\tau;\alpha,\tau)d\alpha+

\\

+\int^{\tau}_0\Big(

K_{1}(\tau,t)u_{x}(l,t)-K_{2}(\tau,t)u(l,t)\Big)dt+\gamma_{1}(\tau),

\label{besh:11}

\end{multline}

где

\begin{gather*}

K_{1}(\tau,t)=\int^{l}_{0}\big[\big(\eta(l,t)w(l,t;\alpha,\tau)\big)_{t}

-k(l,t)w(l,t;\alpha,\tau)\big]d\alpha,

\\

K_{2}(\tau,t)=\int^{l}_{0}\Big[\big(\eta(l,t)w_{x}(l,t;\alpha,\tau)\big)_{t}-

k(l,t)w_{x}(l,t;\alpha,\tau)+r(l,t)w(l,t;\alpha,\tau)\Big]d\alpha,

\end{gather*}

\begin{multline*}

\gamma_{1}(\tau) =\int^{l}_{0}\int^{l}_\alpha\Big(d(x, 0)

w(x, 0;\alpha,\tau)u_{0}(x)-\eta(x, 0)w_{x}(x, 0;\alpha,\tau)u'_{0}(x)

\Big)dxd\alpha-

\\

-\int^{\tau}_0 \int^{l}_{0}

\int^l_{\alpha}w(x,t;\alpha,\tau)f(x,t)dxd\alpha dt+

\int^{l}_{0}u_{x}(l, 0)\eta(l, 0)w(l, 0;\alpha,\tau)d\alpha.

\end{multline*}



Учитывая \eqref{besh:11}, из \eqref{besh:2} получим

\begin{multline}

u(0,\tau)-u(l,\tau)\beta(\tau)\eta(l,\tau)\int^{l}_{0}w_{x}(l,\tau;\alpha,\tau)d\alpha

+u_{x}(l,\tau)\beta(\tau)\eta(l,\tau) \times \\ \times\int^{l}_{0}w(l,\tau;\alpha,\tau)d\alpha

+\int^{\tau}_0\Big(

K_{3}(\tau,t)u_{x}(l,t)+K_{4}(\tau,t)u(l,t)\Big)dt=\gamma_{2}(\tau),

\label{besh:12}

\end{multline}

где

\begin{gather*}

K_{3}(\tau,t)=-\beta(\tau)K_{1}(\tau,t),

\quad 

K_{4}(\tau,t)=\beta(\tau)K_{2}(\tau,t)-\rho(\tau,t),

\\

\gamma_{2}(\tau)=\beta(\tau)\gamma_{1}(\tau)-\mu(\tau).

\end{gather*}



Учитывая \eqref{besh:4}, из представления \eqref{besh:10} при $\alpha=0$ получим интегральное уравнение

\begin{multline}

u(0,\tau)-u(l,\tau)\eta(l,\tau)w_{x}(l,\tau;0,\tau)

+u_{x}(l,\tau)\eta(l,\tau)w(l,\tau;0,\tau) +

\\

+\int^{\tau}_0\Big(K_{5}(\tau,t)u_{x}(l,t)+K_{6}(\tau,t)u(l,t)\Big)dt=\gamma_{3}(\tau),

\label{besh:13}

\end{multline}

где

\begin{gather*}

K_{5}(t,\tau)=-\big(\eta(l,t)w(l,t;0,\tau)\big)_{t}+k(l,t)w(l,t;0,\tau),

\\

K_{6}(t,\tau)=\big(\eta(l,t)w_{x}(l,t;0,\tau)\big)_{t}-

k(l,t)w_{x}(l,t;0,\tau)+r(l,t)w(l,t;0,\tau),

\end{gather*}

\begin{multline*}

\gamma_{3}(\tau)=\int^{l}_{0}\Big(d(x, 0)w(x, 0;0,\tau)u_{0}(x)-

\eta(x, 0)w_{x}(x, 0;0,\tau)u'_{0}(x)\Big)dx -

\\

-\int^{\tau}_0

\int^{l}_{0}w(x,t;0,\tau)f(x,t)dxdt+\eta(l, 0)w(l, 0;0,\tau)u'_{0}(l).

\end{multline*}



Учитывая \eqref{besh:3}, из \eqref{besh:12} и \eqref{besh:13} получим систему интегральных уравнений

\begin{multline}

%\begin{array}{l}

u(0,\tau)-u(l,\tau)\beta(\tau)\eta(l,\tau)\int^{l}_{0}w_{x}(l,\tau;\alpha,\tau)d\alpha

+\int^{\tau}_0\Big(K_{4}(\tau,t)u(l,t)\Big)dt=\gamma_{2}(\tau),

\\

u(0,\tau)-u(l,\tau)\eta(l,\tau)w_{x}(l,\tau;0,\tau)

+\int^{\tau}_0\Big(K_{6}(\tau,t)u(l,t)\Big)dt=\gamma_{3}(\tau), ~~~~~~~

%\end{array}

\label{besh:14}

\end{multline}

которая в операторной форме принимает вид

\begin{equation}

A(\tau)\overrightarrow{u}(\tau)+\int^{\tau}_0B(\tau,t)\overrightarrow{u}(t)

dt=\overrightarrow{\gamma}(\tau),

\label{besh:15}

\end{equation}

где

$$

\det|

A(\tau)|=\beta(\tau)\eta(l,\tau)\int^{l}_{0}w_{x}(l,\tau;\alpha,\tau)d\alpha-

\eta(l,\tau)w_{x}(l,\tau;0,\tau).

$$





\noindent

Покажем, что определитель $\det |A(\tau)|\neq 0$. %\vspace{12pt}

%\pagebreak





\begin{lemma}[1]

%\textbf{Л е м м а 1.} 

Функция $w(x,t;\alpha,\tau)$ удовлетворяет

неравенству

$$

w(x,\tau;0,\tau) >0\ \text{ для любого }\ x \in (0,l],\quad

\eta(l,\tau)w_{x}(l,\tau;0,\tau) >1,

$$

если $d(x,t)<0,$ $\eta(x,t)\geq c_{0} >0$ для любых $(x,t) \in D.$

\end{lemma}



\smallskip



\begin{proof} Следуя рассуждениям \cite{besh:bib:Shhanukov}, рассмотрим задачу

\begin{gather}

(\eta(x,\tau)w_{x}(x,\tau;\alpha,\tau))_{x}

+d(x,\tau)w(x,\tau;\alpha,\tau)=0,

\label{besh:16}

\\

w(\alpha,\tau;\alpha,\tau)=0,

\quad 

w_{x}(\alpha,\tau;\alpha,\tau))={\eta^{-1}(\alpha,\tau)}.

\end{gather}



С помощью принципа максимума и принципа Заремба"--~Жиро из \eqref{besh:16}

получаем $w_{x}(0,\tau;0,\tau)>0$ для любого $(x,t) \in [0,l)$. Тогда из

равенства

$$

\eta(l,\tau)w_{x}(l,\tau;0,\tau) =\eta(0,\tau)w_{x}(0,\tau;0,\tau)-

\int^{l}_{0}d(x,\tau) w(l,\tau;0,\tau)dx

$$

имеем, что если $d(x,t)<0, \ \eta(x,t)\geq c_{0} >0$ для любых

$(x,t) \in D$, то



\bigskip

\hspace{4cm}

$

\eta(l,\tau)w_{x}(l,\tau;0,\tau)>1.

$

\end{proof}



\bigskip



На основании доказанной леммы убеждаемся, что если $\beta(\tau)<0$

для любого $\tau \in [0,T]$, то $\det |A(\tau)|\neq 0$. Поэтому

система уравнений \eqref{besh:14} является системой интегральных уравнений

Вольтерра второго рода, которая безусловно разрешима. Таким образом,

находя из интегральных уравнений Вольтерра $u(0,\tau)=f(\tau)$,

$u(l,\tau)=\varphi(\tau)$, где $f(\tau)$, $\varphi(\tau) \in C^1[0,T]$,

задачу \eqref{besh:1}--\eqref{besh:4} редуцируем к первой начально"=краевой задаче,

однозначная разрешимость которой установлена также в работе \cite{besh:bib:Shhanukov}.

Отсюда следуют существование и единственность решения задачи \eqref{besh:1}--\eqref{besh:4}.

Теорема доказана.

\end{proof}



\smallskip



Заметим, что \eqref{besh:3} можно заменить условием

$$

-u_{x}(l,t)=\int^{t}_{0}\rho_{1}(t,\tau)u(l,\tau)d\tau-\mu_{1}(t),

\quad 0 \leq t\leq T,

$$

где $\rho_{1}(t,\tau)$, $\mu_{1}(t)$ "--- функции, непрерывные на

$[0, T]$.%\vspace{12pt}



\smallskip

\Section[n]{Задача B. Существование и единственность решения задачи B}

Рассмотрим теперь нелокальную краевую задачу, когда условие \eqref{besh:3} в

\textbf{задаче А }заменяется условием вида

\begin{equation}

-u_{x}(l,t)=\beta_{1}(t) u(l,t)-\mu_{1}(t), \quad 0 \leq t\leq T,

\tag{\ref{besh:3}$^*$}

\label{besh:3*}

\end{equation}

где $\beta_{1}(t)$, $\mu_{1}(t)$ "--- функции, непрерывные на

$[0,T]$.%\vspace{12pt}



\smallskip



\begin{theorem}[2]Пусть коэффициенты уравнения \eqref{besh:1} и~граничных условий \eqref{besh:2}{\rm ,} \eqref{besh:3*}{\rm ,} \eqref{besh:4} удовлетворяют условиям гладкости \eqref{besh:5}{\rm ,} $d(x,t)<0$ для любых $(x,t) \in D,$ $\beta(t)<0$ и~$\beta_{1}(t)>0$ для любого $t \in [0, T].$ 

Тогда задача \eqref{besh:1}{\rm ,} \eqref{besh:2}{\rm ,} \eqref{besh:3*}{\rm ,} \eqref{besh:4}

имеет единственное регулярное в $\overline{D}$ решение{\rm .}

\end{theorem}



\smallskip



\begin{proof}

С учётом \eqref{besh:3*} уравнения \eqref{besh:12} и \eqref{besh:13} образуют следующую систему

интегральных уравнений:

\begin{equation}

\begin{array}{l}

\displaystyle u(0,\tau)-u(l,\tau)\Big[\beta(\tau)\eta(l,\tau)\int^{l}_{0}

w_{x}(l,\tau;\alpha,\tau)d\alpha +

\\

\hspace{3cm}

\displaystyle 

 + \beta(\tau)\beta_{1}(\tau)\eta(l,\tau)\int^{l}_{0}

w(l,\tau;\alpha,\tau)d\alpha \Big] +

\\

\hspace{6cm}

\displaystyle 

+\int^{\tau}_0\Big(K_{7}(\tau,t)u(l,t)\Big)dt=\gamma_{4}(\tau),

\\

\displaystyle u(0,\tau)-u(l,\tau)\Big[\eta(l,\tau)w_{x}(l,\tau;0,\tau)+

\beta_{1}(\tau)\eta(l,\tau)w(l,\tau;0,\tau)\Big]+

\\

\hspace{6cm}

\displaystyle +\int^{\tau}_0\Big(K_{8}(\tau,t)u(l,t)\Big)dt=\gamma_{5}(\tau),

\end{array}

\label{besh:17}

\end{equation}

где

\begin{gather*}

K_{7}(\tau,t)=K_{4}(\tau,t)-K_{3}(\tau,t)\beta_{1}(\tau),

\\

\gamma_{4}(\tau)=\gamma_{2}(\tau)-\beta(\tau)\eta(l,\tau)\mu_{1}(\tau)\int^{l}_{0}

w(l,\tau;\alpha,\tau)d\alpha-\int^{\tau}_0\Big(K_{3}(\tau,t)\mu_{1}(\tau)\Big)dt,

\\

K_{8}(\tau,t)=K_{6}(\tau,t)-K_{5}(\tau,t)\beta_{1}(\tau),

\\

\gamma_{5}(\tau)=\gamma_{3}(\tau)-\eta(l,\tau)\mu_{1}(\tau)

w(l,\tau;\alpha,\tau)-\int^{\tau}_0\Big(K_{5}(\tau,t)\mu_{1}(\tau)\Big)dt.

\end{gather*}

Систему интегральных уравнений \eqref{besh:17} перепишем в операторной форме

\begin{equation}

A(\tau)\overrightarrow{u}(\tau)+\int^{\tau}_0B(\tau,t)\overrightarrow{u}(t)

dt=\overrightarrow{\gamma}(\tau), 

\label{besh:18}

\end{equation}

где

\begin{multline*}

\det| A(\tau)|=-\eta(l,\tau)w_{x}(l,\tau;0,\tau)-

\beta_{1}(\tau)\eta(l,\tau)w(l,\tau;0,\tau)+

\\

+\beta(\tau)\eta(l,\tau)\int^{l}_{0}

w_{x}(l,\tau;\alpha,\tau)d\alpha +

\beta(\tau)\beta_{1}(\tau)\eta(l,\tau)\int^{l}_{0}

w(l,\tau;\alpha,\tau)d\alpha.

\end{multline*}



На основании леммы 1 при условии, что если $\beta(\tau)<0$, $\beta_{1}(\tau)>0$ для любого $\tau \in [0, T]$, убеждаемся, что определитель $\det |A(\tau)|\neq 0$. 

Поэтому система уравнений \eqref{besh:18} является системой интегральных уравнений Вольтерра второго рода, которая безусловно разрешима. 

Таким образом, находя из интегральных уравнений Вольтерра $u(0,\tau)=f(\tau)$, $u(l,\tau)=\varphi(\tau)$, где $f(\tau)$, $\varphi(\tau) \in C^1[0, T]$, задачу \eqref{besh:1}, \eqref{besh:2}, \eqref{besh:3*}, \eqref{besh:4} редуцируем к первой начально"=краевой задаче, однозначная разрешимость которой установлена в работе \cite{besh:bib:Shhanukov}. 

Отсюда следуют существование и единственность решения задачи \eqref{besh:1}, \eqref{besh:2}, \eqref{besh:3*}, \eqref{besh:4}. \end{proof}



\smallskip



\Section[n]{Задача C. Существование и единственность решения задачи C}

Рассмотрим теперь нелокальную краевую задачу, когда условие \eqref{besh:3} в~\textbf{задаче А} заменяется условием вида

\begin{equation}

-\Pi(l,t)=\beta_{1}(t) u(l,t)-\mu_{1}(t), \quad 0 \leq t\leq T.

\tag{\ref{besh:3}$^{**}$}

\label{besh:3**}

\end{equation}







\begin{theorem}[3] Пусть коэффициенты уравнения \eqref{besh:1} и

граничных условий \eqref{besh:2}{\rm ,} \eqref{besh:3**}{\rm ,} \eqref{besh:4}  удовлетворяют условиям гладкости

\eqref{besh:5}{\rm ,} $d(x,t)<0$ для любого $(x,t) \in D$ и $\beta(t)<0 $ для любых

$t \in [0,T].$ Тогда задача \eqref{besh:1}{\rm ,} \eqref{besh:2}{\rm ,} \eqref{besh:3**}{\rm ,} \eqref{besh:4} имеет единственное

регулярное в $\overline{D}$ решение{\rm .}

\end{theorem}



\smallskip



\begin{proof}

Проинтегрируем \eqref{besh:9} по $\alpha$ от 0 до $l$. Тогда с учётом \eqref{besh:4} получим

\begin{multline}

\int^{l}_{0}u(\alpha,\tau)d\alpha=

u(l,\tau)\eta(l,\tau)\int^{l}_{0}w_{x}(l,\tau;\alpha,\tau)d\alpha-

\\

- \int^{\tau}_0\Big(K_{9}(\tau,t)\Pi(l,t)+

K_{10}(\tau,t)u(l,t)\Big)dt+\gamma_{6}(\tau),

\label{besh:19}

\end{multline}

где

\begin{gather*}

K_{9}(\tau,t)=\int^{l}_{0}w(l,t;\alpha,\tau)d\alpha,

\\

K_{10}(\tau,t)=\int^{l}_{0}\Big(\big(\eta(l,t)w_{x}(l,t;\alpha,\tau)\big)_{t}-

k(l,t)w_{x}(l,t;\alpha,\tau)+r(l,t)w(l,t;\alpha,\tau)\Big)d\alpha,

\end{gather*}

\begin{multline*}

\gamma_{7}(\tau)=\int^{l}_{0}\int^{l}_\alpha\Big(d(x, 0)w(x, 0;\alpha,\tau)u_{0}(x)-

\eta(x, 0)w_{x}(x, 0;\alpha,\tau)u'_{0}(x) \Big)dxd\alpha-

\\

-\int^{\tau}_0\int^{l}_{0}

\int^l_{\alpha}w(x,t;\alpha,\tau)f(x,t)dxd\alpha dt.

\end{multline*}

Учитывая \eqref{besh:2}, \eqref{besh:3**}, из \eqref{besh:19} получим

\begin{multline}

u(0,t)-u(l,\tau)\beta(\tau)\eta(l,\tau)\int^{l}_{0}w_{x}(l,\tau;\alpha,\tau)d\alpha+

\\

+ \int^{\tau}_0\Big(K_{11}(\tau,t)u(l,t)\Big)dt=\gamma_{8}(\tau),

\label{besh:20}

\end{multline}

где

\begin{gather*}

K_{11}(\tau,t)=\beta(\tau)K_{10}(\tau,t)-\beta(\tau)\beta_{1}(\tau)K_{9}(\tau,t)-\rho(\tau,t),

\\

\gamma_{8}(\tau)=\beta(\tau)\gamma_{7}(\tau)-\mu(\tau)-

\int^{\tau}_{0}\beta(\tau)K_{9}(\tau,t)\mu_{1}(t)dt.

\end{gather*}



При $\alpha=0$ из \eqref{besh:9} с учётом \eqref{besh:3**} получаем

\begin{equation}

u(\alpha,\tau) - u(l,\tau)\eta(l,\tau)w_{x}(l,\tau;0,\tau)

+\int^{\tau}_0\Big(K_{12}(\tau,t)u(l,t)\Big)dt=\gamma_{9}(\tau),

\label{besh:21}

\end{equation}

где

\begin{multline*}

K_{12}(\tau,t)=\big(\eta(l,t)w_{x}(l,t;0,\tau)\big)_{t}-

k(l,t)w_{x}(l,t;0,\tau)+

\\

+r(l,t)w(l,t;0,\tau)-\beta_{1}(\tau)w(l,t;0,\tau),

\end{multline*} 

\begin{multline*}

\gamma_{9}(\tau)=\int^{l}_0\Big(d(x, 0)w(x, 0;0,\tau)u_{0}(x)-\eta(x, 0)w_{x}(x, 0;0,\tau)u'_{0}(x)

\Big)dx-

\\

-\int^{\tau}_0 \int^l_{0}w(x,t;0,\tau)f(x,t)dxdt -

\int^{\tau}_{0}w(l,t;0,\tau)\mu_{1}(t)dt.

\end{multline*}







Уравнения \eqref{besh:20} и \eqref{besh:21} образуют систему интегральных уравнений.

Запишем систему в операторном виде

\begin{equation}

A(\tau)\overrightarrow{u}(\tau)+\int^{\tau}_0B(\tau,t)\overrightarrow{u}(t)

dt=\overrightarrow{\gamma}(\tau),

\label{besh:22}

\end{equation}

где

$$

\det| A(\tau)|=\beta(\tau)\eta(l,\tau)\int^{l}_{0}

w_{x}(l,\tau;\alpha,\tau)d\alpha - \eta(l,\tau)w_{x}(l,\tau;0,\tau).

$$



На основании леммы 1 при условии

$\beta(\tau)<0$ для любых $\tau \in [0,T]$ убеждаемся, что

определитель $\det |A(\tau)|\neq 0$. Поэтому система уравнений \eqref{besh:22}

является системой интегральных уравнений Вольтерра второго рода,

которая безусловно разрешима. Таким образом, находя из интегральных

уравнений Вольтерра $u(0,\tau)=f(\tau)$, $u(l,\tau)=\varphi(\tau)$,

где $f(\tau)$, $\varphi(\tau)\in C^1[0,T]$, задачу \eqref{besh:1}, \eqref{besh:2}, \eqref{besh:3**}, \eqref{besh:4}

редуцируем к первой начально"=краевой задаче, однозначная

разрешимость которой установлена в работе \cite{besh:bib:Shhanukov}. Отсюда следуют

существование и единственность решения задачи \eqref{besh:1}, \eqref{besh:2}, \eqref{besh:3**}, \eqref{besh:4}. \end{proof}



\smallskip



Заметим, что \eqref{besh:3} также можно заменить условием

$$

-\Pi(l,t)=\int^{t}_{0}\rho_{1}(t,\tau)u(l,\tau)d\tau-\mu_{1}(t),

\quad 0 \leq t\leq T.

$$



\smallskip



\begin{remark} Если ввести аналог функции Римана $\nu=\nu(x,t; \xi,\tau)$ для уравнения \eqref{besh:1} в области $\Omega$ в форме

\begin{gather*}

M\nu(x,t;\xi,\tau)=-(\eta(x,t) \nu_{x})_{xt}+(k(x,t)\nu_{x})_{x}

-(r(x,t)\nu)_{x}-(d\nu)_{t}- q(x,t)\nu=0,

\\

\nu(\xi,t;\xi,\tau)=0, \quad 

\nu_{x}(\xi,t;\xi,\tau)={\eta^{-1}(\xi,\tau)}

\exp \biggl({\int^{t}_{\tau}\frac{k(\xi,t_{1})}{\eta(\xi,t_{1})}dt_{1}}\biggr),

\\

\nu(x,\tau;\xi,\tau)=\omega_{2}(x,\tau),

\end{gather*}

где $\omega_{2}(x,\tau)$ "--- решение задачи Коши

\begin{gather*}

(\eta(x,\tau)\nu_{x}(x,\tau;\xi,\tau))_{x} +

d(x,\tau)\nu(x,\tau;\xi,\tau)=0,

\\

\nu(\xi,\tau;\xi,\tau)=0,

\quad 

\nu_{x}(\xi,\tau;\xi,\tau)={\eta^{-1}(\xi,\tau)},

\end{gather*}

то имеет место следующее представление:

\begin{multline}

u(\xi,\tau) = u(0,\tau)\eta(0,\tau)\nu_{x}(0,\tau;\xi,\tau) -

\int^{\tau}_0\Big(\eta(0,t)\nu(0,t;\xi,\tau)u_{xt}(0,t)+

\\ +

k(0,t)\nu(0,t;\xi,\tau)u_{x}(0,t)

+u(0,t)

\Big[\big(\eta(0,t)\nu_{x}(0,t;\xi,\tau)\big)_{t}-k(0,t)\nu_{x}(0,t;\xi,\tau)+

\\

+r(0,t)\nu(0,t;\xi,\tau)\Big]\Big)dt+

\int^{\xi}_0\Big(\eta(x, 0)\nu_{x}(x, 0;\xi,\tau)u_{x}(x, 0)-

\\

-d(x, 0)\nu(x, 0;\xi,\tau)u(x, 0)

\Big)dx+\int^{\tau}_0 \int^{\xi}_0\nu(x,t;\xi,\tau)f(x,t)dxdt.

\label{besh:24}

\end{multline}

Существование и единственность аналога функции Римана доказаны в~\cite{besh:bib:Shhanukov}.

\end{remark}



\smallskip







\begin{lemma}[2]Функция $\nu(x,t;\xi,\tau)$ удовлетворяет неравенству

$$

\nu(x,\tau;l,\tau) <0\ \text{ для любого }\ x \in [0, l),\quad

\eta(0,\tau)\nu_{x}(0,\tau;l,\tau) >1\,,

$$

если $d(x,t)<0,$ $\eta(x,t)\geq c_{0} >0$ для любого $(x,t) \in D.$

\end{lemma}



\smallskip



На основании леммы 2 и представления \eqref{besh:24} аналогично  доказываются существование и~единственность регулярных решений задач A, B, C, в~которых условие \eqref{besh:2} заменяется последовательно следующими условиями:

$$

1)\quad u_{x}(0,t)=\beta(t)\int^{l}_{0}u(x,t)dx +

\int^{t}_{0}\rho(t,\tau)u(l,\tau)d\tau-\mu(t), \quad 0 \leq t \leq T,

$$



$\qquad $ при условии, что $\beta(t)>0$;

$$

2)\quad \Pi(0,t)=\beta(t)\int^{l}_{0}u(x,t)dx +

\int^{t}_{0}\rho(t,\tau)u(l,\tau)d\tau-\mu(t), \quad 0 \leq t \leq T.

$$



{\thanks{Работа выполнена при финансовой поддержке Министерства образования и науки Российской Федерации. Регистрационный номер НИР: \No~1.6197.2011.}



%% Благодарности, гранты и прочее



\begin{thebibliography}{11}



\RBibitem{besh:bib:Barenblat} 

\by Г.~И.~Баренблат, Ю.~П.~Желтов, И.~Н.~Кочина

\paper Об основных представлениях теории фильтрации однородных жидкостей в трещиноватых породах 

\jour ПММ

\yr 1960 

\vol 24

\issue 5 

\pages 852--864

\transl англ. пер.:~

\by G.~I.~Barenblatt, Yu.~P.~Zheltov, I.~N.~Kochina

\paper Basic concepts in the theory of seepage of homogeneous liquids in fissured rocks [strata]

\jour J.~Appl. Math. Mech. 

\vol 24

\issue 5

\yr 1960

\pages 1286--1303 

%http://dx.doi.org/10.1016/0021-8928(60)90107-6





\RBibitem{besh:bib:Dzekzer} 

\paper Уравнения  движения подземных вод со свободной поверхностью в многослойных средах 

\by Е.~С.~Дзекцер

\jour Докл. Акад. наук СССР 

\yr 1975

\vol 220 

\issue 3 

\pages 540--543

\transl англ. пер.:~

\paper Equation of motion of underground water with a free surface in multilayer media

\by E.~S.~Dzektser

\jour Soviet Physics Doklady

\vol 20

\issue 3

\pages 24

\yr 1975

%	1975SPhD...20...24D



\RBibitem{besh:bib:Rubinshtein} 

\paper  К вопросу о процессе распространения тепла в гетерогенных средах 

\by Л.~И.~Рубинштейн

\jour Изв. Акад. наук СССР, Cер. Геогр.

\yr 1948

\vol 12 

\issue 1 

\pages 27--45

\misctransl

\by L.~I.~Rubinstein  

\paper On the problem of the process of propagation of heat in heterogeneous media

\jour Izv. Akad. Nauk SSSR, Ser. Geogr.

\yr 1948

\vol 12 

\issue 1 

\pages 27--45





\RBibitem{besh:bib:TingCooling} 

\by T.~W.~Ting

\paper A cooling process according to two-temperature theory of heat conduction

\jour J.~Math. Anal. Appl.

\yr 1974

\vol 45 

\issue 1

\pages 23–31





\Bibitem{besh:bib:Hallaire} 

\by M.~Hallaire,  S.~de~Parcevaux, R.~J.~Bouchet, et.~al.

\book L'eau et la production v\'eg\'etale

\publaddr Paris 

\yr 1964

\publ Institut National De La Recherche Agronomique

\totalpages 455 

%DOI: 10.1002/qj.49709038625





\RBibitem{besh:bib:Chudnovsky} 

\by А.~Ф.~Чудновский

\book Теплофизика почв 

\publaddr М. 

\publ Наука 

\yr 1976 

\totalpages 352

\misctransl

\by A.~F.~Chudnovsky 

\yr 1976

\book Thermophysics of the soil 

\publaddr Moscow

\publ Nauka

\yr 1976 

\totalpages 352





\Bibitem{besh:bib:Colton} 

\by D.~Colton

\paper Pseudoparabolic equation in one space variable

\jour J.~Diff. Eq.

\yr 1972 

\vol 12 

\issue 3

\pages 559--565





\Bibitem{besh:bib:Colton2} 

\by D.~Colton

\paper Integral operators and the first initial-boundary value problems for pseudoparabolic equations with analytic coefficients

\jour J.~Diff. Eq.

\yr 1973 

\vol 13 

\issue 3

\pages 506--522





\RBibitem{besh:bib:Shhanukov} 

\by М.~X.~Шхануков

\paper О некоторых краевых задачах для уравнения третьего порядка, возникающих при моделировании фильтрации жидкости в пористых средах

\jour Диффер. уравн.

\yr 1982

\vol 18  

\issue 4 

\pages 689--699

\misctransl

\by M.~Kh.~Shkhanukov

\paper On some boundary-value problems for a third-order equation arising when modelling fluid filtration in porous media

\jour Differ. Uravn.

\vol 18

\issue 4

\pages 689–699 

\yr 1982



\RBibitem{besh:bib:Chudnovsky2} 

\by А.~Ф.~Чудновский

\paper Некоторые коррективы в постановке и решении задач тепло- и влагопереноса в почве 

\jour  Сб. трудов АФИ

\yr 1969 

\issue 23 

\pages 41--54

\misctransl

\by A.~F.~Chudnovsky 

\paper Some adjustments in the formulation and solution of problems of heat and moisture transfer in the soil

\jour  Sb. Trudov AFI

\yr 1969 

\issue 23 

\pages 41--54



\RBibitem{besh:bib:Kozhanov} 

\by А.~И.~Кожанов

\paper Об одной нелокальной краевой задаче с переменными коэффициентами для уравнений теплопроводности и Аллера 

\jour Диффер. уравн.

\yr 2004

\vol 40 

\issue 6

\pages 763--774

\transl англ. пер.:~

\by A.~I.~Kozhanov

\paper On a nonlocal boundary value problem with variable coefficients for the heat equation and the Aller equation

\jour Differ. Equ.

\vol 40

\issue 6

\pages 815--826 

\yr 2004







\end{thebibliography}



\noindent

\hskip6.5cm \thedate \par\noindent

\hskip6.5cm\therevdate





\makeengtitle



\newpage

%

% \tableofengcontents

%

%\end{document}

